\documentclass[11pt]{article}

\usepackage[T1]{fontenc}
\usepackage[utf8]{inputenc}
\usepackage{lmodern}
\usepackage{microtype}
\usepackage{geometry}
\usepackage{amsmath,amssymb,amsthm,mathtools}
\usepackage{enumitem}
\usepackage{hyperref}
\usepackage[nameinlink,noabbrev]{cleveref}

\geometry{margin=1in}
\hypersetup{
  hidelinks,
  pdftitle={Selected results from Graph Edge Colouring: Vizing's Theorem and Goldberg's Conjecture},
  pdfauthor={Notes prepared by ChatGPT}
}
\setlist[enumerate]{label=(\roman*),ref=(\roman*)}
\setlength{\parindent}{0pt}
\setlength{\parskip}{0.55em}
\allowdisplaybreaks

\newtheorem{theorem}{Theorem}[section]
\newtheorem{lemma}[theorem]{Lemma}
\newtheorem{proposition}[theorem]{Proposition}
\newtheorem{corollary}[theorem]{Corollary}
\theoremstyle{definition}
\newtheorem{definition}[theorem]{Definition}
\newtheorem{remark}[theorem]{Remark}

\newcommand{\N}{\mathbb N}
\newcommand{\calC}{\mathcal C}
\newcommand{\E}{E}
\newcommand{\V}{V}
\newcommand{\abs}[1]{\lvert #1\rvert}
\newcommand{\set}[1]{\{#1\}}
\newcommand{\ceil}[1]{\left\lceil #1\right\rceil}
\newcommand{\missing}[1]{\overline{#1}}
\newcommand{\boundary}{\partial}
\newcommand{\colours}{[k]}
\newcommand{\phiT}[1]{\varphi\langle #1\rangle}
\newcommand{\GammaG}{\Gamma(G)}

\title{\textbf{Selected results from}\\[0.4em]
  \emph{Graph Edge Colouring: Vizing's Theorem and Goldberg's Conjecture}}
\author{Notes prepared by ChatGPT}
\date{}

\begin{document}
\maketitle

\begin{abstract}
These notes give a self-contained formulation and proof of the material needed from
Chapter~1 through Theorem~1.4, together with the elementary-set theorem for
Tashkinov trees from Section~5.1 of Stiebitz, Scheide, Toft, and Favrholdt.
The exposition is newly written and is not a verbatim reproduction of the book.
Notation is chosen to agree as closely as possible with the accompanying
formalization manuscript.  The proof of the Tashkinov-tree theorem is expanded
at the delicate terminal-path step; McDonald's thesis is used as an independent
verification source for that reduction.
\end{abstract}

\tableofcontents

\section{Scope and notation}

All graphs in these notes are finite loopless multigraphs.  Parallel edges are
allowed.  For a multigraph $G=(V,E)$ and a vertex set $X\subseteq V$, let
$E_G(X)$ be the set of edges with both ends in $X$, and let
$\boundary_G(X)$ be the set of edges with exactly one end in $X$.  We suppress
the subscript $G$ when the ambient graph is clear.  The degree of $v$ is denoted
$d_G(v)$, and $\Delta(G)=\max_{v\in V}d_G(v)$.

A \emph{$k$-edge-colouring} of $G$ is a map
$\varphi:E(G)\to [k]=\{1,\ldots,k\}$ such that incident edges receive distinct
colours.  The set of all such colourings is denoted $\calC^k(G)$, and the
\emph{chromatic index} is
\[
  \chi'(G)=\min\{k:\calC^k(G)\neq\varnothing\}.
\]
Every colour class of a proper edge-colouring is a matching.

For an odd vertex set $U\subseteq V$ with $\abs U\ge 3$, define its density by
\[
  \gamma_G(U)=\frac{2\abs{E_G(U)}}{\abs U-1},
\]
and put
\[
  \Gamma(G)=\max\left\{\gamma_G(U):
  U\subseteq V,\ \abs U\ge 3,\ \abs U\text{ odd}\right\}.
\]
If no such set exists, the maximum may be taken as $0$; this convention will not
matter below.

\section{Density, elementary multigraphs, and critical cores}

\begin{proposition}[The two elementary lower bounds]\label{prop:lower-bounds}
For every loopless multigraph $G$,
\[
  \chi'(G)\ge \Delta(G)
  \qquad\text{and}\qquad
  \chi'(G)\ge \ceil{\Gamma(G)}.
\]
\end{proposition}

\begin{proof}
At a vertex of degree $\Delta(G)$, all incident edges must receive distinct
colours, proving the first inequality.

Let $U\subseteq V(G)$ be odd, with $\abs U\ge3$, and suppose that $G$ has a
proper $k$-edge-colouring.  A matching contains at most
$(\abs U-1)/2$ edges of $E_G(U)$.  Since the $k$ colour classes partition
$E_G(U)$,
\[
  \abs{E_G(U)}\le k\frac{\abs U-1}{2},
  \qquad\text{so}\qquad
  k\ge \frac{2\abs{E_G(U)}}{\abs U-1}.
\]
Taking the maximum over all eligible $U$ and then using integrality gives the
second inequality.
\end{proof}

\begin{definition}[Elementary multigraph]
A multigraph $G$ is \emph{elementary} if
\[
  \chi'(G)=\ceil{\Gamma(G)}.
\]
\end{definition}

\begin{definition}[Critical multigraph]
A multigraph $G$ is \emph{critical} if every proper subgraph $H$ of $G$ satisfies
$\chi'(H)<\chi'(G)$.  It is \emph{$k$-critical} if it is critical and
$\chi'(G)=k+1$.
\end{definition}

\begin{lemma}[Basic facts about critical graphs]\label{lem:critical-facts}
Let $G$ be $k$-critical.
\begin{enumerate}
  \item $G$ has no isolated vertex.
  \item For every edge $e\in E(G)$, the graph $G-e$ has a $k$-edge-colouring.
\end{enumerate}
\end{lemma}

\begin{proof}
If $v$ were isolated, then $G-v$ would be a proper subgraph with exactly the
same edge-colourings as $G$, contrary to criticality.  For an edge $e$,
criticality gives $\chi'(G-e)\le k$.  A colouring with fewer than $k$ colours
may be regarded as a $k$-edge-colouring by adjoining unused colours.
\end{proof}

\begin{proposition}[Critical-subgraph reduction]\label{prop:critical-core}
Let $G$ be a multigraph, and let $H\subseteq G$ be a critical subgraph with
$\chi'(H)=\chi'(G)$.  If $H$ is elementary, then $G$ is elementary.
\end{proposition}

\begin{proof}
By \cref{prop:lower-bounds} and monotonicity of $\Gamma$ under taking
supergraphs,
\[
  \ceil{\Gamma(G)}
  \le \chi'(G)
  =\chi'(H)
  =\ceil{\Gamma(H)}
  \le\ceil{\Gamma(G)}.
\]
Thus all inequalities are equalities, and
$\chi'(G)=\ceil{\Gamma(G)}$.
\end{proof}

\begin{remark}
Every multigraph contains a critical subgraph with the same chromatic index:
repeatedly delete an edge or vertex whenever doing so leaves the chromatic index
unchanged.  Finiteness guarantees termination.
\end{remark}

\begin{proposition}[Critical elementary graphs; book Proposition~1.3]
\label{prop:critical-elementary}
Let $G$ be a critical elementary multigraph with
$\chi'(G)=k+1$ and $k\ge\Delta(G)+1$.  Put $n=\abs{V(G)}$ and
$m=\abs{E(G)}$.  Then:
\begin{enumerate}
  \item $n$ is odd, the whole vertex set witnesses the density bound, and
  \[
    \chi'(G)=\ceil{\Gamma(G)}
    =\ceil{\frac{2m}{n-1}}.
  \]
  \item For every edge $e\in E(G)$ and every
  $\varphi\in\calC^k(G-e)$, each colour in $[k]$ misses exactly one vertex
  of $G$.  In particular, $V(G)$ is $\varphi$-elementary.
  \item
  \[
    m=k\frac{n-1}{2}+1.
  \]
\end{enumerate}
\end{proposition}

\begin{proof}
Since $G$ is elementary, $\ceil{\Gamma(G)}=k+1$, so $\Gamma(G)>k$.
Choose an odd set $U\subseteq V(G)$ such that
\[
  2\abs{E_G(U)}>k(\abs U-1).
  \tag{2.1}\label{eq:critical-dense-set}
\]
We claim that every edge of $G$ has both ends in $U$.  If an edge $f$ did not,
criticality would give a $k$-edge-colouring of $G-f$.  All edges of $E_G(U)$
would still be present, but its $k$ colour classes would be matchings and hence
would contain at most $k(\abs U-1)/2$ edges of $E_G(U)$ in total, contrary to
\eqref{eq:critical-dense-set}.  Thus every edge lies in $E_G(U)$.  A vertex
outside $U$ would consequently be isolated, contrary to
\cref{lem:critical-facts}; hence $U=V(G)$ and $n$ is odd.  In particular,
\[
  m>k\frac{n-1}{2}.
  \tag{2.2}\label{eq:critical-edge-lower}
\]

Fix $e\in E(G)$ and $\varphi\in\calC^k(G-e)$.  Each colour class is a
matching on the odd set $V(G)$ and therefore has at most $(n-1)/2$ edges.
Consequently
\[
  m-1=\abs{E(G-e)}\le k\frac{n-1}{2}.
\]
On the other hand, integrality and \eqref{eq:critical-edge-lower} give the
reverse inequality.  Equality follows.  Every one of the $k$ colour classes
must therefore have exactly $(n-1)/2$ edges, so it saturates all but one
vertex.  This proves the second assertion and also
$m=k(n-1)/2+1$.

Finally,
\[
  \frac{2m}{n-1}=k+\frac{2}{n-1}>k,
\]
so $\ceil{\Gamma(G)}\ge k+1$.  The density lower bound gives
$\ceil{\Gamma(G)}\le\chi'(G)=k+1$.  Thus equality holds and the whole vertex
set witnesses it, proving the first assertion.
\end{proof}

\section{Partial colourings and the equivalence theorem}

Fix an edge $e\in E(G)$ and a colouring
$\varphi\in\calC^k(G-e)$.  For $v\in V(G)$, let
\[
  \varphi(v)=\{\varphi(f): f\in E(G-e)\text{ is incident with }v\},
  \qquad
  \missing\varphi(v)=[k]\setminus\varphi(v).
\]
For $X\subseteq V(G)$, put
\[
  \missing\varphi(X)=\bigcup_{v\in X}\missing\varphi(v).
\]

\begin{definition}[Elementary, closed, and strongly closed sets]
Let $X\subseteq V(G)$.
\begin{enumerate}
  \item $X$ is \emph{$\varphi$-elementary} if
  $\missing\varphi(u)\cap\missing\varphi(v)=\varnothing$ for all distinct
  $u,v\in X$.
  \item $X$ is \emph{$\varphi$-closed} if no colour in
  $\missing\varphi(X)$ occurs on an edge of $\boundary_G(X)$.
  \item $X$ is \emph{$\varphi$-strongly closed} if it is $\varphi$-closed and
  distinct edges of $\boundary_G(X)$ receive distinct colours.
\end{enumerate}
\end{definition}

For a colour $\alpha\in[k]$ and a set $X\subseteq V(G)$, define
\[
  m_\alpha(X)=\abs{\{v\in X:\alpha\in\missing\varphi(v)\}},
  \qquad
  b_\alpha(X)=\abs{\{f\in\boundary_G(X):\varphi(f)=\alpha\}}.
\]
The boundary is the same in $G$ and $G-e$ whenever both ends of $e$ belong to
$X$.

\begin{lemma}[Per-colour counting identity]\label{lem:per-colour-count}
Suppose both ends of $e$ lie in $X$.  For every $\alpha\in[k]$,
\[
  2\abs{E_\alpha(X)}+b_\alpha(X)=\abs X-m_\alpha(X),
  \tag{3.1}\label{eq:per-colour}
\]
where $E_\alpha(X)$ is the set of $\alpha$-coloured edges of $G-e$ with both
ends in $X$.
\end{lemma}

\begin{proof}
Count incidences between vertices of $X$ and $\alpha$-coloured edges.  Exactly
$\abs X-m_\alpha(X)$ vertices of $X$ are incident with an $\alpha$-coloured
edge.  Each internal $\alpha$-edge contributes two incidences, and each
boundary $\alpha$-edge contributes one.
\end{proof}

\begin{theorem}[Book Theorem 1.4; equivalence used as \texttt{egraph2}]
\label{thm:egraph2}\label{egraph2}
Let $G=(V,E)$ be critical with $\chi'(G)=k+1$, where
$k\ge\Delta(G)+1$.  The following statements are equivalent.
\begin{enumerate}
  \item\label{item:egraph-elementary} $G$ is elementary.
  \item\label{item:egraph-every} For every $e\in E$ and every
  $\varphi\in\calC^k(G-e)$, the set $V$ is $\varphi$-elementary.
  \item\label{item:egraph-some} There exist $e\in E$ and
  $\varphi\in\calC^k(G-e)$ for which $V$ is $\varphi$-elementary.
  \item\label{item:egraph-closed} There exist $e\in E$,
  $\varphi\in\calC^k(G-e)$, and $X\subseteq V$ containing both ends of $e$
  such that $X$ is $\varphi$-elementary and $\varphi$-strongly closed.
\end{enumerate}
\end{theorem}

\begin{proof}
\emph{\cref{item:egraph-elementary}$\Rightarrow$\cref{item:egraph-every}.}
This is exactly \cref{prop:critical-elementary}(ii).

\emph{\cref{item:egraph-every}$\Rightarrow$\cref{item:egraph-some}.}
Choose any edge $e$.  By criticality, $\calC^k(G-e)\neq\varnothing$.

\emph{\cref{item:egraph-some}$\Rightarrow$\cref{item:egraph-closed}.}
Take $X=V$.  Its boundary is empty, so it is strongly closed.

\emph{\cref{item:egraph-closed}$\Rightarrow$\cref{item:egraph-elementary}.}
Let $x$ be an end of $e$.  In $G-e$,
\[
  \abs{\missing\varphi(x)}
  =k-d_{G-e}(x)
  \ge k-(\Delta(G)-1)
  \ge 2.
\]
Choose $\alpha\in\missing\varphi(x)$.  Since $X$ is elementary,
$m_\alpha(X)=1$, and since $X$ is closed, $b_\alpha(X)=0$.  The identity
\eqref{eq:per-colour} gives
\[
  2\abs{E_\alpha(X)}=\abs X-1.
\]
Thus $\abs X$ is odd.  Since $X$ contains the two distinct ends of $e$, it has
odd order at least three.

Put $q=\abs{\missing\varphi(X)}$.  Elementarity gives
\[
  q=\sum_{v\in X}\abs{\missing\varphi(v)}.
\]
Summing degrees in $G-e$ over $X$ yields
\[
  q
  =k\abs X-\sum_{v\in X}d_{G-e}(v)
  =k\abs X-2\abs{E_{G-e}(X)}-\abs{\boundary_G(X)}.
  \tag{3.4}\label{eq:missing-sum}
\]
Strong closure has two consequences: no boundary edge uses a colour in
$\missing\varphi(X)$, and the boundary-edge colours are pairwise distinct.
Hence
\[
  \abs{\boundary_G(X)}\le k-q.
\]
Substituting this inequality into \eqref{eq:missing-sum} gives
\[
  2\abs{E_{G-e}(X)}
  =k\abs X-\abs{\boundary_G(X)}-q
  \ge k(\abs X-1).
\]
Because $e$ has both ends in $X$,
\[
  2\abs{E_G(X)}=2\abs{E_{G-e}(X)}+2>k(\abs X-1).
\]
Therefore $\Gamma(G)>k$, so
$\ceil{\Gamma(G)}\ge k+1=\chi'(G)$.  The reverse inequality follows from
\cref{prop:lower-bounds}.  Thus $G$ is elementary.
\end{proof}

\section{Tashkinov's method}

\subsection{Trees, missing colours, and Kempe interchanges}

\begin{definition}[$k$-triple]
A \emph{$k$-triple} is a triple $(G,e,\varphi)$ such that $G$ is $k$-critical,
$k\ge\Delta(G)+1$, $e\in E(G)$, and
$\varphi\in\calC^k(G-e)$.
\end{definition}

A \emph{tree-sequence} is a sequence
\[
  T=(y_0,e_1,y_1,e_2,y_2,\ldots,e_n,y_n)
  \tag{4.1}\label{eq:tree-sequence}
\]
with distinct vertices $y_0,\ldots,y_n$ and distinct edges
$e_1,\ldots,e_n$, where $e_1=e$ joins $y_0$ and $y_1$, and for every
$i\ge2$, the edge $e_i$ joins $y_i$ to some $y_j$ with $j<i$.  Its vertex set
is $V(T)=\{y_0,\ldots,y_n\}$.  The prefix ending at $y_j$ is denoted
\[
  T_j=(y_0,e_1,y_1,\ldots,e_j,y_j).
\]
The edge $e_1$ is uncoloured under $\varphi$; all other tree edges are coloured.
Write
\[
  \phiT{T_j}=\{\varphi(e_i):2\le i\le j\}
\]
for the set of colours used on the coloured edges of $T_j$.

\begin{definition}[Tashkinov tree]
The tree-sequence $T$ is a \emph{Tashkinov tree with respect to
$(G,e,\varphi)$} if, for every $i\ge2$, the colour $\varphi(e_i)$ is missing at
some earlier vertex: there exists $j<i$ such that
$\varphi(e_i)\in\missing\varphi(y_j)$.
\end{definition}

For two colours $\alpha,\beta$, the subgraph of $G-e$ induced by the
$\alpha$- and $\beta$-coloured edges has maximum degree at most two.  Its
components are paths, even cycles, or isolated vertices.  Interchanging
$\alpha$ and $\beta$ on one component preserves properness; this operation is
an \emph{$(\alpha,\beta)$-Kempe interchange}.

Two elementary observations will be used repeatedly.  In a $k$-triple,
\[
  \abs{\missing\varphi(y_0)},\abs{\missing\varphi(y_1)}\ge2,
  \qquad
  \abs{\missing\varphi(v)}\ge1\quad(v\notin\{y_0,y_1\}),
  \tag{4.2}\label{eq:missing-lower}
\]
because $k\ge\Delta(G)+1$ and the edge $e$ is absent at its two ends.
Also, $G$ has no $k$-edge-colouring, since $\chi'(G)=k+1$.

\subsection{Minimal bad trees}

Call a pair $(\psi,S)$ \emph{bad} if
$\psi\in\calC^k(G-e)$, $S$ is a Tashkinov tree with respect to
$(G,e,\psi)$, and $V(S)$ is not $\psi$-elementary.

For a tree-sequence $S=(z_0,f_1,z_1,\ldots,f_m,z_m)$, define its
\emph{path index} $\rho(S)$ to be the least $r\in\{1,\ldots,m\}$ such that
\[
  z_r,f_{r+1},z_{r+1},\ldots,f_m,z_m
\]
is a graph path in the displayed order.  Equivalently, for every $i$ with
$r\le i<m$, the edge $f_{i+1}$ joins $z_i$ and $z_{i+1}$.  The choice $r=m$
is always valid.  Order bad pairs lexicographically by
\[
  (\rho(S),m).
\]

Assume, for a contradiction, that a bad pair exists.  Choose one of minimum
rank and denote it $(\varphi,T)$, with $T$ as in \eqref{eq:tree-sequence}.  Put
$q=\rho(T)$.  Every proper prefix $T_j$ is elementary: otherwise
$(\varphi,T_j)$ would be a bad pair of smaller rank.  The same conclusion
continues to hold after any recolouring that preserves $T$ as a bad
Tashkinov tree of the same rank.

\begin{lemma}[Unused missing colours]\label{lem:four-unused}
Let $T_j$ be a proper prefix of a minimum-rank bad pair.  At least four colours
are missing on $V(T_j)$ and are not used on the coloured tree edges
$e_2,\ldots,e_j$.
\end{lemma}

\begin{proof}
By \eqref{eq:missing-lower}, the total number of missing-colour incidences on
$V(T_j)$ is at least
\[
  2+2+(j-1)=j+3.
\]
Since $T_j$ is elementary, these incidences have distinct colours.  At most
$j-1$ colours occur on the $j-1$ coloured edges $e_2,\ldots,e_j$.
Consequently at least $(j+3)-(j-1)=4$ missing colours are unused on $T_j$.
\end{proof}

\begin{lemma}[Controlled interchange]\label{lem:controlled-swap}
Let $(\varphi,T)$ be a minimum-rank bad pair.  Suppose $i<j<n$, and let
distinct colours $\alpha,\beta$ satisfy
\[
  \alpha\in\missing\varphi(y_i),
  \qquad
  \beta\in\missing\varphi(y_j),
  \qquad
  \alpha,\beta\notin\phiT{T_j}.
\]
Then the $(\alpha,\beta)$-component containing $y_j$ is a path whose other end
is $y_i$.  After interchanging $\alpha$ and $\beta$ on that component,
$T$ remains a bad Tashkinov tree of the same rank; in the new colouring,
$\beta$ is missing at $y_i$ and $\alpha$ is missing at $y_j$.
\end{lemma}

\begin{proof}
Let $Q$ be the $(\alpha,\beta)$-component containing $y_j$.  Since $y_j$
misses $\beta$, it has degree at most one in this two-coloured subgraph and is
an end of $Q$.  Likewise, if $y_i\in V(Q)$, then $y_i$ is an end because it
misses $\alpha$.

Suppose first that $Q$ does not have $y_i$ as its other end.  Then $Q$ avoids
$y_i$.  Interchanging the two colours on $Q$ makes $\alpha$ missing at $y_j$
while leaving it missing at $y_i$.  No edge of $T_j$ has colour $\alpha$ or
$\beta$, so the displayed tree-sequence and all witness conditions for $T_j$
are unchanged.  Thus $T_j$ becomes a non-elementary Tashkinov tree, contrary
to the minimality of $(\varphi,T)$.

Hence $Q$ joins $y_j$ to $y_i$.  Perform the interchange.  The missing statuses
of $\alpha$ and $\beta$ are exchanged at $y_i$ and $y_j$ and are unchanged at
every other vertex of $T$.  The tree edges through $T_j$ are unchanged.  For a
later tree edge $e_h$ with $h>j$, the Tashkinov condition is preserved as
follows.  A witness for a colour outside $\{\alpha,\beta\}$ is unaffected.  If
an old witness for $\alpha$ or $\beta$ is lost at one of $y_i,y_j$, the other
one supplies a witness, and both precede $e_h$.  If $e_h$ itself lies on $Q$
and changes from $\alpha$ to $\beta$, then $y_i$ is an earlier vertex missing
$\beta$ after the interchange; if it changes from $\beta$ to $\alpha$, then
$y_j$ is an earlier vertex missing $\alpha$.  Therefore $T$ remains a
Tashkinov tree.

For each of the two colours, the number of vertices of $T$ missing that colour
is unchanged, and every other missing-colour multiplicity is unchanged.
Thus $V(T)$ remains non-elementary.  The sequence and hence its rank are
unchanged.
\end{proof}

\begin{remark}\label{rem:third-endpoint}
A two-coloured component that joins two specified vertices missing one of its
colours cannot contain a third vertex that also misses one of those colours:
all three would have degree at most one in the component.  We shall use this
without further comment to see that a repeated missing colour at $y_n$ survives
a controlled interchange whose two ends are earlier vertices.
\end{remark}

\begin{lemma}[The minimum bad tree is not a path]\label{lem:not-path}
For the minimum-rank bad pair above, $q>2$.
\end{lemma}

\begin{proof}
If $n=1$, non-elementarity says that some colour is missing at both ends
$y_0,y_1$ of $e_1=e$.  Assigning that colour to $e$ would produce a
$k$-edge-colouring of $G$, a contradiction.  Hence $n\ge2$.

If $q\le2$, then, after possibly interchanging the names of $y_0$ and $y_1$,
the entire displayed sequence is a path.  Indeed, this is immediate for
$q=1$; for $q=2$, the edge $e_2$ joins $y_2$ to one of $y_0,y_1$, and the
initial two vertices may be ordered so that it joins $y_1$ to $y_2$.

Because $T_{n-1}$ is elementary while $T$ is not, some colour $\beta$ is
missing at $y_n$ and at an earlier vertex $y_i$.  Suppose $i<n-1$, and choose
$\gamma\in\missing\varphi(y_{i+1})$.  Both $\beta$ and $\gamma$ are unused on
$T_{i+1}$.  For $\gamma$, any occurrence on a tree edge of this prefix would
have an earlier missing-colour witness, contradicting the elementarity of
$T_{i+1}$.  For $\beta$, elementarity forces the witness for any such
occurrence to be $y_i$; the only possible edge is then $e_{i+1}$, but that edge
is incident with $y_i$, where $\beta$ is missing.  Apply
\cref{lem:controlled-swap} with the colours $\beta,\gamma$.  It moves the
missing colour $\beta$ from $y_i$ to $y_{i+1}$ and, by
\cref{rem:third-endpoint}, leaves $\beta$ missing at $y_n$.

Repeating, we may assume that $\beta$ is missing at both $y_{n-1}$ and $y_n$.
The last edge $e_n$ joins these two vertices, so recolour $e_n$ with $\beta$.
Let $\delta$ be its former colour.  The Tashkinov condition for $e_n$ supplies
an earlier vertex $y_j$ missing $\delta$; because $e_n$ was incident with
$y_{n-1}$, we have $j\le n-2$.  After the recolouring, $\delta$ is also missing
at $y_{n-1}$.  Hence $T_{n-1}$ is a shorter bad Tashkinov path, a contradiction.
\end{proof}

\subsection{Tashkinov's theorem}

\begin{theorem}[Book Theorem 5.1; Tashkinov's theorem]
\label{thm:tashkinov}\label{TashTree}
Let $(G,e,\varphi)$ be a $k$-triple, and let $T$ be a Tashkinov tree with
respect to $(G,e,\varphi)$.  Then $V(T)$ is $\varphi$-elementary.
\end{theorem}

\begin{proof}
Assume otherwise and retain the minimum-rank bad pair and notation introduced
above.  By \cref{lem:not-path}, $q>2$.  Since $T_{n-1}$ is elementary, every
repeated missing colour on $T$ occurs at $y_n$ and at some earlier vertex.
We distinguish three exhaustive cases.

\medskip
\noindent\textbf{Case 1: $q=n$.}
Then the final edge $e_n$ does not join $y_{n-1}$ to $y_n$.  We first put the
colouring into three convenient normal forms.  Each operation below preserves
$T$ as a minimum-rank bad pair, so all prefix-elementarity and unused-colour
conclusions remain available.

\smallskip
\emph{(A) There is a colour $\beta$ missing at $y_n$ and at some
$y_j$ with $j\le n-2$, and $\beta$ is unused on all tree edges.}

Initially choose a repeated missing colour $\beta$ at $y_j,y_n$ with $j<n$.
By \cref{lem:four-unused} applied to $T_{n-2}$, at least four missing colours
are unused there.  The two additional edges $e_{n-1},e_n$ use at most two
colours, so at least two colours missing on $T_{n-2}$ are unused on the whole
of $T$.  Choose one of them, say
$\alpha\in\missing\varphi(y_m)$ with $m\le n-2$ and $\alpha\ne\beta$.

If $j=n-1$, then $\beta$ is unused on $T_{n-1}$: an occurrence there would
have an earlier missing-$\beta$ witness, contradicting the elementarity of
$T_{n-1}$.  It is not used on $e_n$, which is incident with the
$\beta$-missing vertex $y_n$.  Thus both $\alpha$ and $\beta$ are unused on
all tree edges.  A controlled interchange between $y_m$ and $y_{n-1}$ moves
the repeated missing colour $\beta$ to $y_m,y_n$, establishing (A).

Now suppose $j\le n-2$.  If $\beta$ is already unused on $T$, there is nothing
to prove.  Otherwise, let $Q$ be the maximal $(\alpha,\beta)$-chain starting
at $y_n$.  If $Q$ meets $V(T_{n-2})$, traverse it from $y_n$ and let $y_h$ be
its first vertex in that prefix.  Append the segment from $y_h$ to $y_n$ to
$T_{n-2}$.  Its internal vertices lie outside the retained prefix, and its
edges alternate between $\alpha$ and $\beta$.  Both colours are missing in
$T_{n-2}$, at $y_m$ and $y_j$, respectively, so every appended edge has a
valid earlier witness.  The resulting Tashkinov tree is bad because
$\beta$ is still missing at $y_j,y_n$.  Its path index is at most $n-1$, which
is smaller than $q=n$, a contradiction.

Hence $Q$ avoids $V(T_{n-2})$.  Every tree edge has an end in
$V(T_{n-2})$: this is clear for $e_1,\ldots,e_{n-1}$, and it holds for $e_n$
because $e_n$ is not incident with $y_{n-1}$.  Therefore $Q$ contains no tree
edge.  Interchange $\alpha$ and $\beta$ on $Q$, and then globally interchange
the names $\alpha$ and $\beta$.  The local interchange leaves every tree-edge
colour unchanged and changes no earlier witness needed by a tree edge; the
global relabelling plainly preserves the Tashkinov condition.  In the final
colouring, $\beta$ is missing at $y_m,y_n$.  Moreover, any tree edge now
coloured $\beta$ would have been coloured $\alpha$ before the two operations,
but $\alpha$ was unused on $T$.  Thus (A) holds.

\smallskip
\emph{(B) If $\gamma=\varphi(e_n)$, then $\gamma$ is present at
$y_{n-1}$.}

If $\gamma$ is missing at $y_{n-1}$, the Tashkinov condition for $e_n$ and the
elementarity of $T_{n-1}$ imply that $y_{n-1}$ is the unique missing-$\gamma$
vertex in that prefix.  Consequently $\gamma$ is unused on $T_{n-1}$.  The
colour $\beta$ from (A) is unused on all of $T$.  Apply a controlled
$(\beta,\gamma)$-interchange between $y_j$ and $y_{n-1}$.  The component cannot
contain $y_n$: the latter misses $\beta$ and is incident with the
$\gamma$-coloured edge $e_n$, so it would be a third end of the component.
Thus $e_n$ keeps colour $\gamma$, while $\gamma$ becomes present at
$y_{n-1}$.  This proves (B).

\smallskip
\emph{(C) Some colour is missing at $y_n$ and at a vertex of
$T_{n-2}$.}

In the current bad pair, choose a repeated missing colour $\beta'$ at
$y_{j'},y_n$.  If $j'\le n-2$, (C) holds.  If $j'=n-1$, then $\beta'$ is
unused on all of $T$: it is unused on $T_{n-1}$ by elementarity and cannot
colour $e_n$, which is incident with $y_n$.  As above, choose a different
colour $\alpha'$ missing on $T_{n-2}$ and unused on the whole tree.  A
controlled interchange between the missing-$\alpha'$ vertex in $T_{n-2}$ and
$y_{n-1}$ moves $\beta'$ into $T_{n-2}$ while leaving it missing at $y_n$.
Neither colour involved is $\gamma$: the colour $\gamma$ is used on $e_n$, and
by (B) it is present at $y_{n-1}$.  Hence (B) remains true, and (C) is
established.

Delete $y_{n-1}$ and $e_{n-1}$ from the displayed sequence, retaining $e_n$
and $y_n$ as the final edge and vertex:
\[
  T'=(y_0,e_1,y_1,\ldots,e_{n-2},y_{n-2},e_n,y_n).
\]
Because $q=n$, the earlier end of $e_n$ lies in $T_{n-2}$, so $T'$ is a
tree-sequence.  By (B), the missing-$\gamma$ witness for the
$\gamma$-coloured edge $e_n$ is not $y_{n-1}$; it therefore remains in
$T_{n-2}$.  Thus $T'$ is a Tashkinov tree.  By (C), it is bad.  Its length is
$n-1$, so its path index is at most $n-1<n=q$, contradicting minimality.

\medskip
\noindent\textbf{Case 2: $q<n$, and a colour is missing at $y_n$ and at
some $y_i$ with $i\ge q$.}

Let $\beta$ be such a colour.  If $i<n-1$, choose
$\gamma\in\missing\varphi(y_{i+1})$.  Since the suffix from $y_q$ to $y_n$ is
a path, $e_{i+1}$ joins $y_i$ to $y_{i+1}$.  Exactly as in the proof of
\cref{lem:not-path}, both $\beta$ and $\gamma$ are unused on $T_{i+1}$.
A controlled interchange moves the missing colour $\beta$ from $y_i$ to
$y_{i+1}$ and leaves it missing at $y_n$.  Repeating, we arrange that
$\beta$ is missing at both $y_{n-1}$ and $y_n$.

Now recolour the last edge $e_n$ with $\beta$.  Let $\delta$ be its former
colour.  Its Tashkinov witness is some $y_j$ with $j\le n-2$, because
$\delta$ was present at the end $y_{n-1}$ of $e_n$.  After the recolouring,
$\delta$ is also missing at $y_{n-1}$.  Therefore $T_{n-1}$ is a bad
Tashkinov tree of smaller rank, a contradiction.

\medskip
\noindent\textbf{Case 3: $q<n$, and no colour missing at $y_n$ is missing at
any $y_i$ with $i\ge q$.}

Since $T$ is bad and $T_{n-1}$ is elementary, some colour $\beta$ is missing at
$y_n$ and at a vertex $y_i$ with $i<q$.  We may first arrange that
$i\le q-2$.  Indeed, if $i=q-1$, apply \cref{lem:four-unused} to $T_{q-2}$.
After accounting for the single additional edge $e_{q-1}$, at least three
colours missing on $T_{q-2}$ are unused on $T_{q-1}$.  Choose one
$\alpha\ne\beta$, missing at some $y_m$ with $m\le q-2$.  The colour $\beta$
is also unused on $T_{q-1}$, by the elementarity of that prefix.  A
controlled interchange between $y_m$ and $y_{q-1}$ moves the repeated missing
colour $\beta$ to $y_m,y_n$.

Choose a colour $\gamma\in\missing\varphi(y_q)$.  It is unused on $T_q$, for
otherwise a tree edge of that prefix would have an earlier missing-$\gamma$
witness, contradicting elementarity.  Also $\gamma\ne\beta$, since this case
assumes that no colour missing at $y_n$ is missing at $y_q$.

By \cref{lem:four-unused} applied to $T_{q-2}$, after accounting for the two
additional edges $e_{q-1},e_q$, at least two colours missing on $T_{q-2}$ are
unused on $T_q$.  Choose one of them, say
$\varepsilon\in\missing\varphi(y_{m'})$ with $m'\le q-2$, so that
$\varepsilon\ne\gamma$.

If $\varepsilon$ is missing at $y_n$, a controlled
$(\varepsilon,\gamma)$-interchange between $y_{m'}$ and $y_q$ moves the
missing colour $\varepsilon$ to $y_q$ and leaves it missing at $y_n$.  This
reduces the situation to Case~2.

Suppose instead that $\varepsilon$ is present at $y_n$.  Then
$\varepsilon\ne\beta$.  Let $R$ be the maximal
$(\varepsilon,\beta)$-chain starting at $y_n$.  If $R$ meets
$V(T_{q-1})$, traverse it from $y_n$ and let $y_h$ be its first vertex in that
prefix.  If $h\le q-2$, append the segment from $y_h$ to $y_n$ to
$T_{q-2}$; if $h=q-1$, append it to $T_{q-1}$.  In the first subcase, both
$\varepsilon$ and $\beta$ are missing in $T_{q-2}$; in the second, they are
missing in $T_{q-1}$.  Thus every appended edge has an earlier missing-colour
witness.  The resulting tree is bad because $\beta$ is missing at the retained
vertex $y_i$ and at $y_n$.  Its path index is at most $q-1$, contradicting the
minimality of $q$.

It remains that $R$ avoids $V(T_{q-1})$.  Interchange $\varepsilon$ and
$\beta$ on $R$.  No edge of $T_q$ lies on $R$, because every such edge has an
end in $V(T_{q-1})$.  The tree remains Tashkinov: all witness conditions
through $T_q$ are unchanged, and every later edge whose colour or witness is
affected has an earlier witness for $\beta$ at $y_i$ or for $\varepsilon$ at
$y_{m'}$.  After the interchange, $\varepsilon$ is missing at both
$y_{m'}$ and $y_n$.  The colours $\varepsilon$ and $\gamma$ are still unused
on $T_q$, and $\gamma$ is still missing at $y_q$.  A controlled
$(\varepsilon,\gamma)$-interchange between $y_{m'}$ and $y_q$ now makes
$\varepsilon$ missing at $y_q$ while leaving it missing at $y_n$.  Again we
are in Case~2, which has already been shown impossible.

All three cases contradict the choice of the minimum-rank bad pair.  Therefore
no bad pair exists, and every Tashkinov tree in a $k$-triple is elementary.
\end{proof}

\section{Concordance with the formalization target}

The statements above match the two imported results as follows.
\begin{itemize}[leftmargin=2em]
  \item Book Theorem~1.4 (\Cref{thm:egraph2}) is the four-way equivalence
  referred to by the formalization label \texttt{egraph2}; the theorem carries both labels
  \texttt{thm:egraph2} and \texttt{egraph2}.  Its hypotheses include exactly the
  criticality, chromatic-index, and $k\ge\Delta(G)+1$ assumptions used there.
  \item Book Theorem~5.1 (\Cref{thm:tashkinov}) is the assertion that the
  vertex set of every Tashkinov tree in a $k$-triple is elementary; it carries both labels
  \texttt{thm:tashkinov} and \texttt{TashTree}.
  \item \Cref{prop:critical-core} is the critical-subgraph reduction used to
  pass from an arbitrary counterexample to a critical one without changing
  the chromatic index.
\end{itemize}

\begin{thebibliography}{9}

\bibitem{SSFT2012}
Michael Stiebitz, Diego Scheide, Bjarne Toft, and Lene M. Favrholdt,
\emph{Graph Edge Colouring: Vizing's Theorem and Goldberg's Conjecture},
Wiley, 2012.

\bibitem{McDonald2009}
Jessica McDonald,
\emph{A Survey of Recent Results on the Chromatic Index of Multigraphs},
Ph.D. thesis, University of Waterloo, 2009.

\bibitem{Sebo2024}
Andr\'as Seb\H{o},
``Tashkinov-Trees: an Annotated Proof,''
\emph{Graphs and Combinatorics} \textbf{40} (2024), Article~18,
\href{https://doi.org/10.1007/s00373-023-02712-1}{doi:10.1007/s00373-023-02712-1}.

\end{thebibliography}

\end{document}
